\documentclass[10pt,a4paper]{amsart}
\usepackage[margin=19mm]{geometry}
\usepackage{amsmath,amssymb,mathtools}
\usepackage[scr=rsfs,cal=euler]{mathalfa}
\usepackage{xfrac}
\usepackage[unicode,hidelinks,bookmarks=false]{hyperref}
\newcommand{\ie}{i.e.\ }
\newcommand{\cf}{cf.\ }
\newcommand{\resp}{respectively\ }
\newcommand{\Subsection}{\subsection}
\providecommand{\nobreakdash}{\nobreak\hbox{-}}
\setlength{\emergencystretch}{2em}
\multlinegap0pt
\usepackage{amssymb}
\usepackage{bm}
\usepackage{tikz}
\newtheorem{lem}{Lemma}
\newtheorem{claim}{Claim}
\title{Linear continuous operators with bounded supports}
\author{Vesko Valov}
\date{Source: arXiv:2604.25228v4 (CC BY 4.0)}
\begin{document}
\maketitle

\noindent\emph{Source and licence.} Linear continuous operators with bounded supports. Version arXiv:2604.25228v4 (CC BY 4.0), distributed by arXiv under the Creative Commons Attribution 4.0 International licence, http://creativecommons.org/licenses/by/4.0/, as recorded in the arXiv licence metadata for this exact version and shown on the abstract page https://arxiv.org/abs/2604.25228v4. Full licence notice: \texttt{licenses/arxiv-2604.25228-CC-BY-4.0.txt}.\par\smallskip

\noindent\textit{Editorial scope.} Counted unit: the article's lem lem (source0.tex lines 456--461). The article's complete original proof of it is delivered with it (source0.tex lines 462--527, one \texttt{proof} environment of 66 lines). No mathematics is written, repaired or re-derived: the statement and the proof are the article's own lines, in the article's own notation, and no retained line is substituted or re-flowed.  No further source-local context is needed: every cross-reference of every retained line resolves inside the two delivered spans.  Nothing was omitted from any retained span, so the package carries no omission record.  No retained line contains a citation, so the wrapper carries no bibliography and no bibliography entry.  No retained line contains a figure environment, an image-inclusion command or drawing code, so no image is delivered and the package declares no asset dependency.  Editorial formatting only, in addition: an AMS \texttt{amsart} preamble that restates the article's own declarations and macros used by the retained lines, namely the environments \texttt{lem}, \texttt{claim} on separate counters, together with the standard packages those lines need.  The retained environments print in this document as Lemma 1, Claim 1 in order of appearance, whereas the article prints the same items with the numbering of its own full document.  The complete native source of the article as distributed in the arXiv e-print for this exact version is bundled unchanged as \texttt{sources/199284/source0.tex}.
\begin{lem}
Let $Y$ be a space and $D(Y)\subset C(Y)$. Denote by $D(\overline Y)$ the set of all extensions $\overline g\in C(\beta Y,\overline{\mathbb R})$, $g\in D(Y)$.
Then for every countable 
set $\overline\Phi_0\subset D(\overline Y)$ there is  a countable set $\overline\Phi\subset D(\overline Y)$ containing $\overline\Phi_0$ such that all real-valued elements of $E_{\overline\Phi}=\{\pi_{\overline g}: \overline g\in\overline\Phi\}$ is dense in
$C_p((\triangle\overline\Phi)(\beta Y))$ and the set $\Phi=\{\overline g|Y:\overline g\in\overline\Phi\}$ is admissible. 
\end{lem}

\begin{proof} Observe that $D(\overline Y)=C(\beta Y)$ if $D(Y)=C^*(Y)$.
% and $D(\overline Y)\subset C(\beta Y,\overline{\mathbb R})$ such that $\{\overline g|Y;\overline g\in C(\beta Y)\}= D(Y)$ otherwise.
\begin{claim}
For every countable $\Psi'\subset D(Y)$ there is a countable admissible set $\Psi\subset D(Y)$ containing $\Psi'$.
\end{claim}
We construct by induction an increasing sequence of countable sets $\Psi_n\subset D(Y)$ each containing $\Psi'$ and maps 
$\delta^{n+1}_n:(\triangle\Psi_{n+1})(Y)\to(\triangle\Psi_{n})(Y)$ such that: 
\begin{itemize}
\item Every space $Y_n=(\triangle\Psi_{n})(Y)$ has a countable base $\mathcal B_n$ with $(\delta^{n+1}_n)^{-1}(\mathcal B_n)\subset\mathcal B_{n+1}$;
\item $\Psi_{2n-1}$ is closed under addition, multiplication and multiplication by rational numbers;
\item For every pair $U,V\in\mathcal B_{2n-1}$ with $cl_{Y_n}(V)\subset U$ there exists a  fixed $f_{U,V}\in C^*(Y_{2n-1})$ such that 
$f_{U,V}(V)=1$, $f_{U,V}(Y_{2n-1}\backslash U)=0$ and $f_{U,V}\circ(\triangle\Psi_{2n-1})\in D(Y)$; 
\item $\Psi_{2n}=\Psi_{2n-1}\cup\{f_{U,V}\circ\triangle\Psi_{2n-1}:U,V\in\mathcal B_{2n-1}\}$.
\end{itemize} 
Since $D(Y)$ is an algebra, there is a countable $\Psi_1\subset D(Y)$ containing $\Psi'$ such that $\Psi_1$ is closed under addition, multiplication and multiplication by rational numbers. Then 
$\overline\Psi_1=\{\overline g:g\in\Psi_1\}\subset D(\overline Y)$ and $Y_1=(\triangle\Psi_{1})(Y)$ is a dense subset of 
$\overline Y_1=(\triangle\overline\Psi_{1})(\beta Y)$. We fix a countable base $\widetilde{\mathcal B}_1$ of $\overline Y_1$ and for every $\widetilde U,\widetilde V\in\widetilde{\mathcal B}_1$ with $cl(\widetilde V)\subset\widetilde U$ choose  
a function $\overline f_{\widetilde U,\widetilde V}\in C(\overline Y_1)$ such that  
%$\overline f_{\widetilde U,\widetilde V}\circ\triangle\overline\Psi_{2n-1}\in F(\overline Y)$,
$\overline f_{\widetilde U,\widetilde V}(\widetilde V)=1$ and $\overline f_{\widetilde U,\widetilde V}(\overline Y_1\backslash \widetilde U)=0$.
Denote $\mathcal B_1=\{\widetilde U\cap Y_1:\widetilde U\in\widetilde{\mathcal B}_1\}$ and $f_{U,V}=\overline f_{\widetilde U,\widetilde V}|Y_1$.

Observe that $\overline f_{U,V}\circ(\triangle\overline\Psi_{1})\in C(\beta Y)$ and $\big(\overline f_{U,V}\circ(\triangle\overline\Psi_{1})\big)|Y=f_{U,V}\circ(\triangle\Psi_{1}(Y))$. Then
each $f_{U,V}\circ(\triangle\Psi_{1}(Y))$ belongs to $D(Y)$. Next, define $\overline\Psi_2=\overline\Psi_1\cup\{\overline f_{\widetilde U,\widetilde V}\circ(\triangle\overline\Psi_{1}):\widetilde U,\widetilde V\in\widetilde{\mathcal B}_1\}$. Because $\overline \Psi_1\subset\overline\Psi_2$ there is natural map $\overline\eta^2_1:\overline Y_2=(\triangle\overline\Psi_2)(\beta Y)\to \overline Y_1$. Choose a base $\widetilde{\mathcal B}_2$ for $\overline Y_2$ which is closed under finite intersections and containing all 
$(\overline\eta^2_1)^{-1}(\widetilde U)$, $\widetilde U\in\widetilde{\mathcal B}_1$. Let $\Psi_2=\overline\Psi_2|Y_2$, $\mathcal B_2=\{\widetilde U\cap Y_2:\widetilde U\in\widetilde{\mathcal B}_2\}$ and $\delta^2_1=\eta^2_1|Y_2$.
Now, it is clear how to complete the construction of the sets $\Psi_n$. 

Let show that the countable set 
$\Psi=\bigcup_{n=1}^\infty\Psi_n$ is admissible. Since $\{\Psi_n\}$ is increasing and each $\Psi_{2n-1}$ is closed under addition, multiplication and 
multiplication by rationals, so is $\Psi$. That implies the set $E_\Psi=\{\pi_g:g\in\Psi\}$ is also closed under addition, multiplication and multiplication by rationals. Hence, we need to show that for every $y\in Y_\Psi=(\triangle\Psi)(Y)$ and its neighborhood $U\subset Y_\Psi$ there exists $g\in\Psi$ with
$\pi_g(y)=1$ and $\pi_g(Y_\Psi\backslash U)=0$. To this end, observe that $Y_{\overline\Psi}=(\triangle\overline\Psi)(\beta Y)$ is the limit space of the inverse sequence
$S_\Psi=\{\overline Y_n;\eta^{n+1}_n)\}$ and, according to our construction, the base of $Y_{\overline\Psi}$ consists of all sets of the form $\eta_n^{-1}(\widetilde U_n)$ with $\widetilde U_n\in\widetilde{\mathcal B}_n$, $n\geq 1$, where 
$\eta_n:Y_{\overline\Psi}\to\overline Y_n$ are the projections in $S_\Psi$. Therefore, we can assume that $U=\widetilde U\cap Y_\Psi$ and $\widetilde U=\eta_n^{-1}(\widetilde U_n)$ with $\widetilde U_n\in\widetilde{\mathcal B}_n$ for some $n$. Passing to a bigger integer, we can suppose that $n=2k-1$. Choose a neighborhood $V_n\in\mathcal B_n$ of $y_n=\eta_n(y)$ with $cl_{Y_n}(V_n)\subset U_n=\widetilde U_n\cap Y_n$. Then the function $f_{U_n,V_n}$ satisfies the equalities $f_{U_n,V_n}(V_n)=1$ and $f_{U_n,V_n}(Y_n\backslash U_n)=0$. Consequently, $g=f_{U_n,V_n}\circ(\triangle\Psi_{n})\in\Psi_{n+1}\subset\Psi$ and $\pi_g=f_{U_n,V_n}\circ\eta_n$. 
Hence, $\pi_g(y)=1$ and $\pi_g(Y_\Psi\backslash U)=0$. This complete the proof of Claim 6. $\diamondsuit$



Next step of our proof is to construct an increasing sequence of countable sets $\overline\Phi_n\subset C(\beta Y,\overline{\mathbb R})$ containing $\overline\Phi_0$ such that:
\begin{itemize}
\item The real-valued elements of each $E_{2n-1}=\{\pi_{\overline g}: \overline g\in\overline\Phi_{2n-1}\}$ is dense in
$C_p((\triangle\overline\Phi_{2n-1})(\beta Y))$;
\item Every $\Phi_{2n}=\{\overline g|Y:\overline g\in\overline\Phi_{2n}\}$ is an admissible set.
\end{itemize}
Let show how to construct the sets $\overline\Phi_{n}$. We choose a set $\overline\Phi_{1}$ according to Lemma 3.3. Then 
$\Phi_{1}=\{\overline g|Y:\overline g\in\overline\Phi_{1}\}$ is a countable subset of $D(Y)$. Hence, by Claim 6, there is a countable admissible set
$\Phi_2\subset D(Y)$ containing $\Phi_{1}$. Let $\overline\Phi_2=\{\overline g:g\in\Phi_2\}\subset C(\beta Y,\overline{\mathbb R})$ be the extension of $\Phi_2$. Next, apply Lemma 3.3 to find a countable set $\overline\Phi_3\subset C(\beta Y,\overline{\mathbb R})$ containing $\overline\Phi_2$ such that
the real-valued elements of $E_{3}=\{\pi_{\overline g}: \overline g\in\overline\Phi_{3}\}$ is dense in
$C_p((\triangle\overline\Phi_{3})(\beta Y))$ and denote $\Phi_3=\{\overline g|Y:\overline g\in\overline\Phi_{3}\}$. In this way, applying either Lemma 3.3 or Claim 6, we construct the sequence $\{\overline\Phi_n\}$. Let $\overline\Phi=\bigcup_{n=1}^\infty\overline\Phi_n$ and $\Phi=\bigcup_{n=1}^\infty\Phi_n$. Since $\Phi$ is the union of the increasing sequence $\{\Phi_{2n}\}$ consisting of admissible sets, $\Phi$ is also admissible, see $(3.4)$. 

So, it remains to show that the following claim:
\begin{claim}
The real-valued elements of $E_{\overline\Phi}=\{\pi_{\overline g}: \overline g\in\overline\Phi\}$ is dense in the space $C_p((\triangle\overline\Phi)(\beta Y))$. 
\end{claim}
To this end let $\overline Y_n=(\triangle\overline\Phi_n)(\beta Y)$ and $\overline Y_0=(\triangle\overline\Phi)(\beta Y)$. For every $n$ there are natural maps $\eta^{2n+1}_{2n-1}:\overline Y_{2n+1}\to \overline Y_{2n-1}$ since $\overline\Phi_{2n-1}\subset\overline\Phi_{2n+1}$. Then 
$\overline Y_0$ is the limit space of the inverse sequence $S_\Phi=\{\overline Y_{2n+1};\eta^{2n+1}_{2n-1}\}$.
Every projection $\eta_{2n-1}:\overline Y_0\to\overline Y_{2n-1}$ induces a continuous map 
$\eta_{2n-1}^*:C_p(\overline Y_{n-1})\to C_p(\overline Y_0)$ defined by $\eta_{2n-1}^*(h)=h\circ\eta_{2n-1}$. 
Because 
%the set of real-valued functions from  
%$E_{2n-1}=\{\pi_{\overline g}:\overline g\in\overline\Phi_{2n-1}\}$ is dense in $C_p(\overline Y_{2n-1})$ and 
$E_{\overline\Phi}=\bigcup_n\eta_{2n-1}^*(E_{2n-1})$ (recall that $\overline\Phi=\bigcup_{n=1}^\infty\overline\Phi_{2n-1}$), it suffices to show that  
the set of real-valued functions from $\bigcup_{n=1}^\infty\eta_{2n-1}^*(E_{2n-1})$ is dense in $C_p(\overline Y_0)$. So, let 
$O=\{g\in C_p(\overline Y_0):|g(\overline y_i)-g_0(\overline y_i)|<\varepsilon_i, i=1,..,k\}$ be a neighborhood of some $g_0\in C_p(\overline Y_0)$, where 
all $\overline y_i$ are different points from $\overline Y_0$. Since the base of $\overline Y_0$ consists of all sets of the form $\eta_{2n-1}^{-1}(U_{2n-1})$ where $U_{2n-1}$ belongs to the base of $\overline Y_{2n-1}$, there is $m$ and different points $y_i\in\overline Y_{2m-1}$ such that $\eta_{2m-1}(\overline y_i)=y_i$ and $\eta_{2m-1}^{-1}(U_i)\subset g_0^{-1}(V_i)$, where $U_i$ are neighborhoods of $y_i$ in $\overline Y_{2m-1}$ and $V_i$ are the open intervals $(g_0(\overline y_i)-\varepsilon_i,g_0(\overline y_i)+\varepsilon_i)$. Then the set $W=\{h\in C_p(\overline Y_{2m-1}):h(y_i)\in V_i, i=1,2,..,k\}$  is open in $C_p(\overline Y_{2m-1})$. Since the set of real-valued functions from $E_{2n-1}$ is dense in $C_p(\overline Y_{2n-1})$, there is a real-valued function
$h_0\in E_{2m-1}$ with $h_0\in W$. Then $h_0\circ\eta_{2m-1}$ is a real-valued function from $\bigcup_{n=1}^\infty\eta_{2n-1}^*(E_{2n-1})\cap O$. So, the set of all real-valued functions from $E_{\overline\Phi}$ is dense in $C_p(\overline Y_0)$.
\end{proof}

\end{document}
